\section{Referee-round positive closure: full gauge, Orlicz duality, and the Gaussian tangent}
\label{sec:rr2-b3}

\subsection{Primal and dual spaces}

Let \(\mathcal X_T\) be the state space of B2.  Collision densities are placed
in the Orlicz space associated with
\(\ell(q)=q\log q-q+1\), and collision sources in its exponential dual.
Density test potentials belong to the weighted Sobolev cylinder space
\(\mathcal Y\) for which
\(\partial_t p+v\cdot\nabla_xp\) and \(\Delta_\omega p\) are bounded by the
velocity weight.

For a triple \((q,p,\psi)\), define
\[
 \mathcal H_f(q,p,\psi)
 =
 \langle q,f\rangle+\langle v\cdot\nabla p,f\rangle
 +\int(e^{\Delta p+\psi}-1)\,dA_f .
\]

\begin{theorem}[Closed path-space Fenchel duality]
\label{thm:rr2-b3-duality}
For every finite-energy density path,
\[
 \sup_{q,p,\psi}
 \left\{
 \langle q,f\rangle+
 \langle\psi,\Gamma\rangle-
 \int_0^T\mathcal H_f(q,p,\psi)\,dt
 \right\}
\]
equals \(\int\ell(d\Gamma/dA_f)dA_f\) if the weak balance and symmetry
constraints hold, and is \(+\infty\) otherwise.  The supremum may be taken
over smooth compactly supported test triples and is stable under weighted
Orlicz approximation.
\end{theorem}

\begin{proof}
Integration by parts in time and space converts the \(q,p\) terms into the
weak balance defect.  If the defect is nonzero, scaling a separating test
potential makes the supremum infinite.  On the constraint set, these terms
cancel and the remaining supremum is the integral of the pointwise conjugate
of \(e^z-1\).  Rockafellar's integral conjugacy applies because the
exponential/entropy pair is a decomposable Orlicz dual pair.  Velocity
truncation and mollification approximate admissible dual variables, while
entropy coercivity gives uniform integrability and lower semicontinuity.
\end{proof}

\subsection{The complete representation gauge}

Let \(r\in\mathcal Y\).  The transformation
\[
 (q,p,\psi)\longmapsto
 \left(
 q-(\partial_t+v\cdot\nabla_x)r,\,
 p+r,\,
 \psi-\Delta_\omega r
 \right)
\]
leaves the action pairing invariant, up to the explicitly retained endpoint
term.  Collision invariants are contained in this gauge.

Define the closed operator
\[
 \mathfrak G r=
 \big(-(\partial_t+v\cdot\nabla)r,\ r,\ -\Delta r\big)
\]
and the cotangent space
\[
 \mathfrak T^*=
 (\mathcal Y_q\times\mathcal Y_p\times\mathcal Y_\psi)/
 \overline{\operatorname{Ran}\mathfrak G}.
\]

\begin{theorem}[Gauge-complete cotangent identification]
\label{thm:rr2-b3-gauge}
Assume the macro-constraint derivative has closed range and satisfies the
Robinson interior condition at a unique finite-rate minimizer.  Then a
Lagrange multiplier exists and its class in \(\mathfrak T^*\) is unique.
The collision relation determines the gauge-invariant quantity
\[
 \Delta p+\psi=\log\frac{d\Gamma}{dA_f}.
\]
There is a unique minimal-norm representative orthogonal to
\(\overline{\operatorname{Ran}\mathfrak G}\).
\end{theorem}

\begin{proof}
The closed-range and interior hypotheses give the Banach-space multiplier
theorem.  If two triples represent the same multiplier, their difference
annihilates every feasible tangent direction.  The adjoint closed-range
theorem identifies this annihilator with
\(\overline{\operatorname{Ran}\mathfrak G}\), proving uniqueness of the
class.  Strict convexity in \(\Gamma\) gives the displayed collision
relation.  Hilbertizing the weighted dual space by the reference quadratic
form gives an orthogonal decomposition and hence a unique minimal-norm
representative.
\end{proof}

\subsection{First and second derivatives of the deterministic pressure}

Let \(Q_\epsilon(\Theta)=\mu_\epsilon^{-1}
\log\mathbb E e^{\mu_\epsilon\Theta}\).  For macroscopic cylinder directions
\(F,G\),
\[
 D^2Q_\epsilon(\Theta)[F,G]
 =
 \mu_\epsilon\operatorname{Cov}_{\Theta,\epsilon}(F,G).
\]

\begin{theorem}[Convergence of pressure derivatives]
\label{thm:rr2-b3-derivatives}
On a compact regular source ball, the marked cluster expansion converges
normally in a complex neighborhood.  Consequently
\[
 DQ_\epsilon(\Theta)\to DQ(\Theta),\qquad
 D^2Q_\epsilon(\Theta)\to D^2Q(\Theta)
\]
uniformly, and the limiting Hessian is the covariance of the
\(\sqrt{\mu_\epsilon}\)-fluctuation field, including the constrained initial
covariance from B1.
\end{theorem}

\begin{proof}
Normal convergence of the analytic polymer series permits termwise
differentiation and Cauchy estimates uniform in \(\epsilon\).  Uniform
convergence on a slightly larger complex ball implies convergence of the
first two derivatives.  The finite-volume differentiation identity supplies
the factor \(\mu_\epsilon\).  At time zero, B1 identifies the quadratic form
as the reference covariance projected onto the tangent space of the
constraints; the dynamic terms are the second marked cumulants of the
cluster series.
\end{proof}

\begin{theorem}[Fluctuating Boltzmann tangent]
\label{thm:rr2-b3-fluctuation}
Let
\[
 \zeta^\epsilon=\sqrt{\mu_\epsilon}(f^\epsilon-f).
\]
Under the regular biased law associated with a smooth cotangent class,
\(\zeta^\epsilon\) converges in the weighted distribution space to the unique
Gaussian solution
\[
 d\zeta_t=\mathcal L_{f_t}\zeta_t\,dt+dW_t^{f},
\]
with initial covariance equal to the constrained B1 covariance and collision
noise covariance
\[
 \mathbb E\langle W^f(\phi),W^f(\chi)\rangle_t
 =
 \int_0^t\!\int
 (\Delta\phi)(\Delta\chi)\,dA_{f_s}\,ds .
\]
\end{theorem}

\begin{proof}
The convergence of all second derivatives in
Theorem~\ref{thm:rr2-b3-derivatives} identifies finite-dimensional Gaussian
limits.  Third and higher centered cumulants of
\(\zeta^\epsilon\) vanish by the Cauchy bounds and the
\(\mu_\epsilon^{-1/2}\) scaling.  Exponential compact containment from B2
gives tightness in the weighted distribution space.  Differentiating the
biased kinetic equation identifies the drift \(\mathcal L_f\); the collision
second derivative gives the displayed martingale bracket.  The constrained
initial LDP supplies the initial Gaussian field.  Uniqueness of the linear
martingale problem identifies the limit.
\end{proof}
